% WHAT AN OVERFULL BOX IS GIVEN.
%
% A horizontal box that is too wide by more than \hfuzz is given a RULE at
% its right, as wide as \overfullrule, and that rule goes INSIDE the box:
%
%   \overfullrule=5pt \hfuzz=0pt
%   \setbox0=\hbox to 0pt{A}
%   Overfull \hbox (2.0pt too wide) detected at line 6
%   \rip A|
%
%   \hbox(7.0+1.0)x0.0
%   .\rip A
%   .\rule(*+*)x5.0
%
% It is put there AFTER the box has been packed, so its own width is no
% part of what was measured -- the box above is still `2.0pt too wide',
% which is the A alone.
%
% Three things get none: a box where \overfullrule is nothing; one that is
% reported only because \hbadness is under a hundred, without being over by
% more than \hfuzz; and a VERTICAL box, however overfull.
%
% HSTeX gave none at all, so every line of trip's that did not fit was one
% node short -- twenty of them in the second pass, and the short display of
% each lost its `|'.
\catcode`\{=1 \catcode`\}=2
\tracingonline=1 \showboxbreadth=10 \showboxdepth=10
\font\rip=trip \rip
\message{[A overfull, with an overfull rule]}
\overfullrule=5pt \hfuzz=0pt
\setbox0=\hbox to 0pt{A}\showbox0
\message{[B overfull, with no rule asked for]}
\overfullrule=0pt
\setbox0=\hbox to 0pt{A}\showbox0
\message{[C overfull within \noexpand\hfuzz, so reported only by badness]}
\overfullrule=5pt \hfuzz=3pt \hbadness=99
\setbox0=\hbox to 0pt{A}\showbox0
\message{[D underfull, which draws no rule]}
\hfuzz=0pt \hbadness=0
\setbox0=\hbox to 20pt{A}\showbox0
\message{[E a vertical box, which draws none either]}
\overfullrule=5pt \vfuzz=0pt \vbadness=0
\setbox0=\vbox to 0pt{\hbox{}\kern10pt}\showbox0
\message{[done]}
\end
